\section{空间信息在真实降阶预测中能省下多少报警？}
\paragraph{数据与比较口径。}
沿用原导热研究的 24 个训练参数、12 个调参参数及 $r=8$ 初始投影基。三个网格 $24^2,48^2,96^2$ 测试同一序列的 256 个参数。旧序列早已公开，因而在这里作为配对复算数据。阈值固定为 $\tau=10^{-3}$（无量纲温升）。程序对该组预测分别评估五种界，没有依据该次复算选择块数。另在 $48^2$ 网格生成 128 个不同的新工况，分布与旧流一致，固定随机种子为 20261121，独立重置初始基。两条序列均未用于本次调参。

同一个残差可揭示证书本身是否变紧。我们在固定八维基和旧自适应模型的回退前预测两种轨迹上，为\emph{同一预测}逐一换用五种界。另让每一种判据从同一个初始基出发，独立决定是否完整求解，并只用本策略实际支付的回退解更新基。后一种比较才衡量策略部署后的回退次数。比较所用的全阶参照解在决定之后核对误差，不提供给判据。原流各网格重复计时三次，重复时清空学习状态。

\begin{table}[ht]\centering\small
\caption{阈值相同，固定旧预测进行配对诊断，或让各策略实际更新自己的基。三元组按 $24^2/48^2/96^2$ 排列；每格均为 256 次查询。}
\begin{tabular}{lrr}\toprule
上界 & 在旧轨迹上报警 & 各自部署后实际回退\\\midrule
原解析标量 & 15 / 15 / 17 & 15 / 15 / 17\\
最优统一倍数 & 15 / 15 / 17 & 15 / 15 / 17\\
按位置分块 & 6 / 11 / 13 & 12 / 14 / 14\\
按位置和正负分开 & 6 / 10 / 13 & 12 / 14 / 14\\
双端点有符号包络 & 6 / 10 / 13 & 12 / 14 / 14\\\bottomrule\end{tabular}
\end{table}

\paragraph{符号确实改变了一次决定。}
在 $48^2$ 网格的旧轨迹上，第 11 次预测的四边散热系数从左、右到下、上依次是 $1.8110$、$1.0467$、$3.4027$ 和 $3.7954$（取近似值）。其最大场误差仅约 $1.73\times10^{-6}$。解析标量界为 $0.00462038$，只分位置时为 $0.00102159$，均会报警；保留正负后降为 $0.00092555$，这次预测可以直接返回。表中 $47\to29$ 是固定旧预测得到的报警数；各自部署后为 $47\to40$。提前增广基会改变其后的残差，故两个问题给出不同的计数。

\paragraph{误差覆盖及仍然保守的情形。}
在原流的独立部署中，五种方法最终返回的结果均达标。双端点的 40 次回退前预测也都已达标。转到固定八维基，三个网格分别有 6、6、7 次实际超标，五种界却仍对每格的全部 256 次预测报警；界值变小未必改变决策。双端点相对原界的逐例比值中位数分别为 $0.196,0.187,0.183$。仅保留符号相对只用位置的中位比值约为 $0.68$，而双端点在这批实际预测上与单端符号界相同。正负互杂的块可能无法提供非零的端点修正。在补充的 128 个新工况中，原方法回退 15 次，分块符号和双端点各回退 13 次；三种策略均纠正了一次真实超标，未观察到漏报。
